\section{Engineering Decisions and Rationale}
\label{sec:decisions}

\subsection{Why beam 2 rather than beam 1?}
Beam~1 is faster (3.46\x{} against 2.20\x{} on KDE4) but loses 2.15 chrF overall and about six points for Chinese and
Japanese. Beam~2 keeps every language within 1.4 points of beam~5 for more than twice the speed.

\subsection{Why not larger batches or cross-file batching?}
Batch layout changed speed by at most 3\% while doubling peak memory, and a 4,096-token attempt did not complete on the
8~GB machine. Batches are therefore kept at 1,024 tokens and per file, which also preserves per-file telemetry.

\subsection{Why square markers for placeholders?}
They survived translation in 99.15\% of cases against 88.78\% and 80.66\% for the alternatives, because the model copies
citation-style markers such as \code{[1]} unchanged.

\subsection{Why keep untranslatable strings in English?}
Forcing the model to avoid \code{<unk>} produced fluent but wrong words. An English string in a localized application is
visible and can be fixed by a translator; a wrong word is silent.

\subsection{Why a hash as the memory key?}
SQL Server limits a key to 900 bytes (434 characters). A fixed 32-byte SHA-256 hash removes the limit while keeping
lookups indexed.

\subsection{Why record the engine version with each translation?}
It makes reuse safe across model and setting changes without manual cleanup: older entries simply stop matching.

\subsection{Why base punctuation rules on counts?}
The first, guessed rules made Japanese reuse worse than fresh translation. Counting how every reference ends turned the
result from a loss into a gain.

\subsection{Why test the memory with commit-like versions?}
A memory only pays off across runs. Versions with genuinely new strings, run in sequence with the memory kept, measure
the savings a CI server would actually see, rather than a seeded cache rate.

\subsection{Why a stopped ready runner?}
A stopped machine keeps its disk (the built image and the model) for \$3.65 per month and starts in seconds, saving about
two minutes per CI run over a fresh machine.
